\section{Cutoff changes and cached witnesses}\label{sec:cutoff}

During search the incumbent improves and the cutoff $U$ decreases. Since
$T_{U'}(B)\subseteq T_U(B)$ for $U'\le U$, a lower cutoff can make
further tightening possible, and certificates computed for the old cutoff
may expire. This section answers three questions about stored primal
information. At which cutoff does a certificate expire? Which
current-round witnesses can be produced for a new cutoff without solving
a new support problem? How do stored points behave when the box changes?

Primal witnesses give upper bounds on the possible benefit of tightening.
Stored dual solutions give the opposite information, a guaranteed
reduction, for a relaxation whose constraint matrix does not change.
Section~\ref{sec:constraints} treats dual information and combines both
kinds in a bracket for one frozen relaxation. Throughout this section, a
\emph{lifted pool} $\widehat W$ is a finite set of complete lifted points
with stored objective values. When a lifted pool is used on one box $B$,
all of its points lie in $R(B)$; mixing and frontiers also require
$R(B)$ to be convex and the objective $v$ to be affine, as in the
reference family.

\subsection{Cutoff thresholds of a protected box}\label{sec:cutoff-thresholds}

Let $P=[a,b]$ be a protected box with a lifted pool
$\widehat W\subseteq R(P)$ whose projections attain every endpoint of
$P$, and write $q_z=v(z)$ for $z\in\widehat W$. The $2n$ \emph{faces} of
$P$ are $\{x\in P:x_i=a_i\}$ and $\{x\in P:x_i=b_i\}$; when $a_i=b_i$ the
two coincide, which does no harm below. The certificate of
Theorem~\ref{thm:protected} remains valid at every cutoff
$U'\ge U_{\widehat W}=\max_{z\in\widehat W}q_z$, the largest stored value.
This threshold can be larger than necessary, because an expensive point
may be redundant. The exact threshold for reusing part of the pool is
\begin{equation}\label{eq:face-threshold}
 \tau_{\widehat W}(P)=\max_F\ \min\{q_z:z\in\widehat W,\ \pi(z)\in F\},
\end{equation}
where $F$ ranges over the faces of $P$. Each inner minimum is over a
nonempty set because $\widehat W$ attains every face.

\begin{proposition}[Face threshold of a pool]\label{prop:face-threshold}
For a cutoff $U'$, the subpool $\{z\in\widehat W:q_z\le U'\}$ attains every
face of $P$ if and only if $U'\ge\tau_{\widehat W}(P)$. Convex
combinations do not lower this threshold: the projection of
$\conv(\widehat W)\cap\{v\le U'\}$ has box hull $P$ if and only if
$U'\ge\tau_{\widehat W}(P)$.
\end{proposition}
\begin{proof}
The subpool attains a face $F$ exactly when some pool point on $F$ has
value at most $U'$, that is, when the inner minimum
in~\eqref{eq:face-threshold} is at most $U'$. Requiring this for all
faces gives the first claim. For the second, the subpool is contained in
$\conv(\widehat W)\cap\{v\le U'\}$, which proves one direction.
Conversely, let $\bar z=\sum_z\lambda_zz$ be a convex combination with
$v(\bar z)\le U'$ whose projection lies on the face $x_i=b_i$. Every
$z\in\widehat W$ has $x_i(z)\le b_i$, and the weighted average equals
$b_i$, so every $z$ with $\lambda_z>0$ lies on that face. Hence
$v(\bar z)=\sum_z\lambda_zq_z$ is at least the inner minimum for that
face. The same argument applies to every face.
\end{proof}

For example, let $R(P)=P=[0,1]^2$ with no auxiliary variables,
$v(x)=x_1+x_2$, and $\widehat W$ the four corners. Then
$U_{\widehat W}=2$ but $\tau_{\widehat W}(P)=1$, because the three
corners $(0,0)$, $(1,0)$ and $(0,1)$ attain all four faces. Mixing can
produce witnesses for a \emph{smaller} box at a lower cutoff, but such a
box needs its own check on its rebuilt relaxation
(Section~\ref{sec:cutoff-mixing}).

The pool threshold depends on the stored points. The following threshold
depends only on the box and the relaxation.

\begin{proposition}[Exact expiry cutoff of a box]\label{prop:box-threshold}
Let $P$ be a nonempty box in a lifted family. For each face $F$ of $P$ let
$\mu_F(P)=\min\{v(z):z\in R(P),\ \pi(z)\in F\}$, with value $+\infty$
when no lift projects to $F$, and let $\tau(P)=\max_F\mu_F(P)$. Then
\[
 T_U(P)=P\quad\Longleftrightarrow\quad U\ge\tau(P).
\]
Moreover, $\tau(P)\le\tau_{\widehat W}(P)$ for every lifted pool
$\widehat W\subseteq R(P)$ that attains every face of $P$.
\end{proposition}
\begin{proof}
Because $R(P)$ is compact and $v$ continuous, $K_U(P)=\pi(R_U(P))$, and
the minimum defining $\mu_F(P)$ is attained when finite. A face $F$ is
therefore attained by a point of $K_U(P)$ exactly when
$\mu_F(P)\le U$. If $K_U(P)$ is nonempty, Proposition~\ref{prop:fixed-box}
turns this face-by-face statement into the claimed equivalence. If
$K_U(P)$ is empty, then $T_U(P)=\emptyset\ne P$, and every face has
$\mu_F(P)>U$. Pool points are feasible for the face problems, which gives
$\tau(P)\le\tau_{\widehat W}(P)$.
\end{proof}

The quantities $\mu_F(P)$ are the finite-box counterparts of the face
values of Corollary~\ref{cor:face-test}. The two thresholds have
different costs and uses. The face threshold $\tau_{\widehat W}(P)$ uses
only the stored coordinates and objective values, without an optimization
solve. When the incumbent value drops below it, the stored certificate
expires. This is an exact event at which reconsidering OBBT is justified;
it does not show that tightening will pay off. The box threshold
$\tau(P)$ requires up to $2n$ face-restricted minimizations of the
objective, comparable to one OBBT round. It gives a positive statement:
if the current box is $P$ and the cutoff $U$ falls below $\tau(P)$, the
next exact round either finds $K_U(P)$ empty or, since some face is no
longer attained and $K_U(P)$ is compact, moves at least one endpoint
strictly. The size of that movement is not bounded below by this
argument; Section~\ref{sec:constraints} gives guaranteed reductions from
dual information. For cutoffs between $\tau(P)$ and
$\tau_{\widehat W}(P)$ the pool is inconclusive. None of these thresholds
is necessary for the existence of a different protected box inside $P$.
Proposition~\ref{prop:threshold-con} is the single-direction counterpart
of $\tau(P)$ for one frozen support problem.

\subsection{Witnesses at a lower cutoff in one relaxation}\label{sec:cutoff-mixing}

A stored current-round witness whose value exceeds a new cutoff can be
moved below it by convex combination with any relaxed point of lower
value.

\begin{proposition}[Mixing to a new cutoff]\label{prop:mixing}
Let $R(B)$ be convex with an affine objective $v$, and let $a,z\in R(B)$
with $v(a)\le U'<v(z)$. Then
\begin{equation}\label{eq:mixing}
 \theta=\frac{U'-v(a)}{v(z)-v(a)}\in[0,1),\qquad
 z_{U'}=a+\theta(z-a)
\end{equation}
satisfies $z_{U'}\in R(B)$ and $v(z_{U'})=U'$.
\end{proposition}
\begin{proof}
The point $z_{U'}$ is a convex combination of two points of the convex
set $R(B)$. Since $v$ is affine,
$v(z_{U'})=v(a)+\theta(v(z)-v(a))=U'$.
\end{proof}

The anchor $a$ needs only relaxation feasibility. A minimizer of $v$ over
$R(B)$, with value $L(B)$, or a lift of the new incumbent are natural
choices. The original coordinates move to
$x(a)+\theta\bigl(x(z)-x(a)\bigr)$, and
Proposition~\ref{prop:current-round} then bounds the frozen round at the
new cutoff without a new support solve. The construction is confined to
one frozen relaxation, and the hull of mixed points is generally not
protected.

\begin{example}[Mixed witnesses at a new cutoff]\label{ex:mixing-hull}
Relax $s=x^2$ on $[-1,1]$ by the square rows $s\ge2|x|-1$ and $s\le1$,
with objective $v=s$.
Mixing the graph points $(\pm1,1)$ with the anchor $(0,0)$ to the cutoff
$U'=1/4$ gives $\theta=1/4$ and the points $(\pm1/4,1/4)$, which are
feasible on $[-1,1]$. On their hull $[-1/4,1/4]$ the secant reads
$s\le1/16$, so both stored lifts fail the rebuilt check. Their projections
nevertheless have feasible rebuilt lifts $(\pm1/4,1/16)$, so this hull is
protected at cutoff $1/4$. On $[-1,1]$ the mixed points bound
the decrease of each endpoint in the frozen round at cutoff $1/4$ by
$3/4$; the exact round gives $[-5/8,5/8]$, a decrease of $3/8$.

A hull of mixed points need not be protected. Mixing the same graph
points with the relaxation minimizer $(0,-1)$ to cutoff $0$ gives
$\theta=1/2$ and points $(\pm1/2,0)$. These attain the supports of the
current cutoff relaxation. On their hull $[-1/2,1/2]$, however,
$\phi(\pm1/2)=1/4>0$, so no feasible lift at cutoff $0$ reaches either
face. The hull is not fixed, and the next round gives $[-1/4,1/4]$.
\end{example}

\subsection{The cutoff frontier of a cached pool}\label{sec:cutoff-frontier}

A nonempty lifted pool $\widehat W=\{z_1,\ldots,z_k\}\subseteq R(B)$ with
values $q_j=v(z_j)$ supplies the current-round witnesses
\[
 \mathcal C_U=\conv(\widehat W)\cap\{v\le U\}\subseteq R_U(B)
\]
at every cutoff $U$. The best ceilings that the pool can give are the
coordinate extrema over $\mathcal C_U$. They require no new support
solve, because they are attained at explicit points. The proof below is
standard extreme-point reasoning; its use here is to generate the
candidates from a retained pool.

\begin{proposition}[Exact frontier of a pool]\label{prop:frontier}
If $\mathcal C_U$ is nonempty, which holds exactly when $\min_jq_j\le U$,
every minimum and maximum of an original or lifted coordinate over
$\mathcal C_U$ is attained at one of the following points:
\begin{enumerate}[label=(\alph*)]
\item a pool point $z_j$ with $q_j\le U$;
\item for a pair with $q_j<U<q_l$, the mixture of $z_j$ and $z_l$ with
  value exactly $U$.
\end{enumerate}
There are at most $k+k(k-1)/2$ such points.
\end{proposition}
\begin{proof}
Let $\Lambda_U=\{\lambda\in\R^k:\lambda\ge0,\ \sum_j\lambda_j=1,\
\sum_j\lambda_jq_j\le U\}$. Because $v$ is affine, the map
$\lambda\mapsto\sum_j\lambda_jz_j$ sends $\Lambda_U$ onto $\mathcal C_U$,
and $\Lambda_U$ is nonempty exactly when some $q_j\le U$. A coordinate is
a linear function of $\lambda$, so its extrema over the polytope
$\Lambda_U$ are attained at vertices. Let $\lambda$ be a vertex with
support $J=\{j:\lambda_j>0\}$. If $|J|\ge2$ and the cutoff inequality is
slack, a nonzero $h$ supported on $J$ with $\sum_jh_j=0$ gives
$\lambda\pm\epsilon h\in\Lambda_U$ for small $\epsilon>0$, which is
impossible at a vertex. If the cutoff is tight and $|J|\ge3$, a nonzero
$h$ supported on $J$ satisfies the two equations $\sum_jh_j=0$ and
$\sum_jh_jq_j=0$, with the same contradiction. If the cutoff is tight,
$J=\{j,l\}$ and $q_j=q_l$, the first equation implies the second, and
again a perturbation exists. Hence either $\lambda$ is a unit vector
$e_j$ with $q_j\le U$, or $J=\{j,l\}$, $q_j\ne q_l$, and the cutoff is
tight. In the latter case two positive weights average $q_j$ and $q_l$ to
$U$, so $U$ lies strictly between them and the image of $\lambda$ is the
stated mixture.
\end{proof}

The proof covers repeated points, equal values, and degenerate hulls. The
proposition optimizes over the convex hull of the pool, not over the
whole relaxation; an empty frontier says only that the pool cannot
supply a witness at that cutoff.

The frontier also describes exactly how the pooled ceilings respond to
the cutoff. Fix a coordinate $i$ and let
$g(U)=\max\{x_i(z):z\in\mathcal C_U\}$ for $U\ge q_{\min}=\min_jq_j$.
The pooled ceiling on the decrease of $u_i$ is $u_i-g(U)$.

\begin{proposition}[Cutoff response of a pool]\label{prop:cutoff-response}
For a nonempty pool, the function $g$ is nondecreasing, concave, and
continuous on $[q_{\min},\infty)$. It is affine between consecutive
distinct values among $q_1,\ldots,q_k$ and constant for
$U\ge\max_jq_j$. The coordinate minimum over $\mathcal C_U$ has the same
properties with \emph{nonincreasing} and \emph{convex} in place of
\emph{nondecreasing} and \emph{concave}.
\end{proposition}
\begin{proof}
The set $\Lambda_U$ grows with $U$, so $g$ is nondecreasing. If $\lambda$
and $\lambda'$ are optimal at $U$ and $U'$, then for $t\in[0,1]$ the
combination $t\lambda+(1-t)\lambda'$ is feasible at $tU+(1-t)U'$, which
gives concavity. A concave function is continuous on the interior of its
interval. At $q_{\min}$, let $\Delta>0$ be the gap to the next larger
pool value; if there is none, $g$ is constant. For
$q_{\min}<U<q_{\min}+\Delta$, every $\lambda\in\Lambda_U$ puts total
weight $\epsilon\le(U-q_{\min})/\Delta$ on points with $q_j>q_{\min}$.
Writing $c=\max_j|x_i(z_j)|$, this gives $g(U)\le g(q_{\min})+2c\epsilon$,
and continuity follows. On an open interval between consecutive distinct
pool values, the sets of eligible points and of straddling pairs in
Proposition~\ref{prop:frontier} do not change. There $g$ is the maximum of
finitely many affine functions of $U$: constants for eligible points, and
for a pair $(j,l)$ the coordinate
$x_i(z_j)+\frac{U-q_j}{q_l-q_j}\bigl(x_i(z_l)-x_i(z_j)\bigr)$. Hence $g$
is convex there; being also concave, it is affine, and by continuity it
is affine on the closed interval. For $U\ge\max_jq_j$, $\Lambda_U$ is the
whole simplex. The statements for the minimum follow by applying these
arguments to $-x_i$.
\end{proof}

Proposition~\ref{prop:cutoff-response} turns a pool into an exact
screening rule for future incumbents. For a tolerance $\delta\ge0$, the
upper direction of $x_i$ is screened at cutoff $U$ when
$g(U)\ge u_i-\delta$. Because $g$ is nondecreasing and continuous, the
set of such cutoffs is empty or an interval $[\sigma_i,\infty)$. The
threshold $\sigma_i$ is found by evaluating $g$ at the distinct pool
values and solving one affine equation. When the incumbent value falls
below $\sigma_i$, the pool can no longer screen this direction, and a
support solve or new witnesses are needed. The computation involves no
optimization problem; at each pool value it evaluates at most
$k+k(k-1)/2$ candidates.

\begin{example}[A frontier and its cutoff response]\label{ex:frontier}
Let $R(B)=B=[0,2]\times[0,3]$ with no auxiliary variables, $v(x)=x_2$,
and the pool $z_1=(0,0)$, $z_2=(2,3)$, $z_3=(1,1)$ with values $0,3,1$.
At $U=2$ the frontier consists of $z_1$, $z_3$, the mixture
$(4/3,2)$ of $z_1$ and $z_2$, and the mixture $(3/2,2)$ of $z_3$ and $z_2$.
The best pooled witness for the upper endpoint of $x_1$ uses the anchor
$z_3$, not the point of lowest value. For this endpoint,
\[
 g(U)=\begin{cases}U,&0\le U\le1,\\ (U+1)/2,&1\le U\le3,\\ 2,&U\ge3.\end{cases}
\]
At $U=2$ the pooled ceiling on the decrease of $u_1=2$ is $1/2$, whereas
the exact decrease is $0$ because $(2,0)\in R_2(B)$. The pool bounds the
round correctly but not tightly. With tolerance $\delta=1/4$, the pool
screens this direction exactly for the cutoffs $U\ge5/2$.
\end{example}

\paragraph{Changed boxes.}
Before a pool is used on a new box, every point must pass the rows of
the relaxation built on that box. For the filtered pool,
Propositions~\ref{prop:frontier} and~\ref{prop:cutoff-response} are again
exact. They are not exact for the old convex hull intersected with the
new rows: mixtures of excluded points can be feasible, and additional
inequalities create coefficient vertices with more than two positive
weights. Under lifted monotonicity, a point feasible on a box $P$ remains
feasible on every box containing $P$, so the witnesses of a protected
box can join the pool at every containing node without a check. Points
computed on a parent box need not be feasible on a child box; in
Example~\ref{ex:halving}, the witness $(r/2,0)$ from $[-r,r]$ fails on
$[-r/2,r/2]$.

\subsection{Continuing after a cutoff decrease}\label{sec:cutoff-continue}

When the cutoff decreases, a solver can either continue tightening from
its current box or restart from the starting box. Continuing yields boxes
contained in the restart iterates. If the current box was obtained by a
finite number of relaxation-sound steps at compatible cutoffs, both
sequences have the same intersection, by monotonicity alone.

\begin{proposition}[Continuing versus restarting]\label{prop:continue}
Let $U'\le U$, let $B_k'=T_{U'}^k(B_0)$, and let
$C_0\subseteq B_0$ be obtained from $B_0$ by $m\ge0$
relaxation-sound steps at cutoffs at least $U'$, for example by OBBT
at the old cutoff $U$. Then $C_k=T_{U'}^k(C_0)$ satisfies
\[
  B_{k+m}'\subseteq C_k\subseteq B_k'\qquad(k\ge0).
\]
In particular, $\bigcap_k C_k=\bigcap_k B_k'$.
\end{proposition}
\begin{proof}
A relaxation-sound step $C\to C'$ at cutoff $U''\ge U'$ gives
$C'\supseteq T_{U''}(C)\supseteq T_{U'}(C)$. Induction and
monotonicity give $C_0\supseteq T_{U'}^m(B_0)$, and hence
$C_k\supseteq T_{U'}^{k+m}(B_0)$. The upper inclusion follows from
$C_0\subseteq B_0$ and monotonicity. The two bounding subsequences have
the same intersection, proving the last statement.
\end{proof}

If the iterates are nonempty and~\eqref{eq:closed-family} holds at the new
cutoff, the common intersection is also the greatest fixed box, by
Proposition~\ref{prop:fixed-limit}. These extra assumptions are needed
for that identification, not for the equality of intersections.

After an incumbent improvement, therefore, no tightening is lost by
resuming from the current box. A box that is protected at the new cutoff,
for example one whose face threshold $\tau_{\widehat W}(P)$ does not
exceed it, still bounds all remaining movement.

\subsection{Conditions for reuse}\label{sec:cutoff-reuse}

Table~\ref{tab:reuse} collects the conditions under which stored primal
information remains valid. The entries restate
Propositions~\ref{prop:current-round}, \ref{prop:face-threshold}
and~\ref{prop:box-threshold}, Theorem~\ref{thm:protected},
Corollary~\ref{cor:original-points}, and the residual certificates of
Section~\ref{sec:residual}.

\begin{table}[htbp]
\centering
\small
\begin{tabular}{@{}>{\raggedright\arraybackslash}p{0.2\linewidth}>{\raggedright\arraybackslash}p{0.23\linewidth}>{\raggedright\arraybackslash}p{0.25\linewidth}>{\raggedright\arraybackslash}p{0.24\linewidth}@{}}
\toprule
Information & Bounds & Valid for & Recheck or expiry\\
\midrule
Pool in $K_U(B)$ & movement in the frozen round on $B$ &
the same box and rows; lower cutoffs through mixing and the frontier &
any change of box or rows\\
\addlinespace
Protected box $P$ with lifted pool $\widehat W$ & all movement of
relaxation-sound steps; relaxation bound through $L(P)$ &
boxes containing $P$; cutoffs at least $\tau_{\widehat W}(P)$ with the
stored pool, or at least $\tau(P)$ with a recomputed pool; the same
family &
box not containing $P$; cutoff below the threshold; new rows violated by
a witness; rounding of fractional faces\\
\addlinespace
Hull of original feasible points $S$ & movement under every reduction
valid for the original problem & boxes containing $S$; cutoffs at least
$\max_{x\in S}f(x)$ & cutoff below $\max_{x\in S}f(x)$; box not
containing $S$\\
\addlinespace
Residual and comparison matrix & total movement of exact Jacobi
iteration & one cutoff and one family, inside the invariant region on
which the matrix is verified & leaving that region; a change of cutoff
or family without a new residual bound and comparison matrix\\
\bottomrule
\end{tabular}
\caption{Conditions for reusing stored primal information. A failed or
expired condition makes the certificate inapplicable; it does not show
that further tightening is useful.}\label{tab:reuse}
\end{table}
